% Profundización y proyecto. Se incluye desde manual_teorico_pedagogico.tex.

\section{Mapas, flujos, señales y sistemas avanzados}

\begin{objectivebox}
Integrar teoría y cómputo para estudiar mapas discretos, flujos caóticos canónicos,
rutas al caos, geometría fractal, datos temporales, retardo, alta dimensión,
dinámica conservativa y transitoria, sincronización, control y memoria fraccionaria.
La meta es formular conclusiones graduadas: qué se derivó, qué se observó, qué
indicador se estimó y qué permanece abierto.
\end{objectivebox}

\subsection{Mapa logístico y diagrama de telaraña}

\paragraph{Ecuación del mapa logístico.}
Un \term{mapa dinámico} actualiza el estado en instantes discretos:
\[
 x_{n+1}=f(x_n;\mu).
\]
La órbita de \(x_0\) es la sucesión
\(\{x_0,f(x_0),f^2(x_0),\ldots\}\). Un punto fijo satisface
\(f(x^*)=x^*\); un punto de periodo primo \(p\) satisface
\(f^p(x)=x\) y no vuelve antes. Para un punto fijo hiperbólico,
\[
 |f'(x^*)|<1\quad\Longrightarrow\quad\text{atractor local},\qquad
 |f'(x^*)|>1\quad\Longrightarrow\quad\text{repulsor local}.
\]
La derivada es un \term{multiplicador por iteración}; las tasas por unidad de tiempo corresponden a los flujos continuos.

\paragraph{Construcción de telaraña.}
En el plano \((x_n,x_{n+1})\), se dibujan \(y=f(x)\) y la diagonal \(y=x\).
Desde \(x_n\), un segmento vertical llega a \(f(x_n)\) y uno horizontal regresa a la
diagonal, donde queda \(x_{n+1}\). La alternancia reproduce exactamente la iteración;
representa la actualización discreta del mapa. Una convergencia alternante indica multiplicador negativo,
mientras que una convergencia sin alternancia suele asociarse con uno positivo.

\paragraph{Ejemplo algebraico resuelto: mapa logístico.}
Para
\[
 f_r(x)=rx(1-x),\qquad 0\leq x\leq1,
\]
los puntos fijos satisfacen
\[
 x=rx(1-x)\quad\Longrightarrow\quad
 x_0^*=0,\qquad x_1^*=1-\frac1r\quad(r\ne0).
\]
Como \(f_r'(x)=r(1-2x)\),
\[
 f_r'(0)=r,\qquad
 f_r'(x_1^*)=2-r.
\]
El origen es estable para \(|r|<1\); el punto no nulo es estable cuando
\(|2-r|<1\), es decir, \(1<r<3\). En \(r=3\) su multiplicador cruza \(-1\), firma
local de la primera duplicación de periodo.

Para \(r=4\) y \(x_0=0.2\),
\[
 x_1=0.64,\qquad x_2=0.9216,\qquad
 x_3=4(0.9216)(0.0784)=0.28901376.
\]
Los tres valores son reproducibles a mano; la interpretación requiere el cálculo de
indicadores además de la apariencia de la gráfica. El exponente finito de una órbita se estima mediante
\[
 \lambda_N=\frac1N\sum_{n=0}^{N-1}
 \log|r(1-2x_n)|.
\]

\begin{activitybox}{Práctica reproducible: cuatro regímenes logísticos}
Seleccione el sistema logístico registrado o defínalo como mapa en
\gui{Crear sistema}. Use \(x_0=0.2\), descarte 500 iteraciones y conserve 500.
Ejecute \(r=2.8\), \(3.2\), \(3.5\) y \(4\). Para los dos primeros dibuje a mano
las primeras veinte etapas de la telaraña; para todos exporte la serie y calcule
\(\lambda_N\). En \gui{Bifurcación}, barra \(r\in[2.5,4]\) con una malla que después
duplique. Conserve el mismo transitorio y número de muestras al comparar.
\end{activitybox}

\begin{figure}[H]
\centering
\includegraphics[width=0.86\linewidth]{logistic_bifurcation.png}
\caption{Diagrama logístico producido para el manual con el software. Cada columna
contiene iterados tardíos para un parámetro; las líneas aparentes son conjuntos de
puntos muestreados, no curvas de soluciones continuas.}
\label{fig:logistica-profundizacion}
\end{figure}

\paragraph{Lectura de la órbita logística.}
Una rama única corresponde a un valor tardío dominante; dos ramas, a una órbita de
periodo dos, si se confirma el orden temporal. Una banda densa puede contener órbitas
caóticas, inestables, preperiódicas y efectos de resolución. Las ventanas claras
dentro de bandas densas motivan una simulación local con mayor transitorio.

\begin{warningbox}
El diagrama de bifurcación debe leerse junto con \(x_0\), transitorio, aritmética, malla y dirección del barrido. La identificación de un atractor único exige
telaraña con muchas líneas tampoco demuestra sensibilidad; para ello se comparan
órbitas vecinas o se estima expansión con control de convergencia.
\end{warningbox}

\lecturaseccion{Steven H. Strogatz, \emph{Nonlinear Dynamics and Chaos with
Applications to Physics, Biology, Chemistry, and Engineering},
2.\textsuperscript{a} ed., secciones 10.1, ``Fixed Points and Cobwebs'', 10.2--10.5,
sobre mapa logístico, ventanas y exponente de Lyapunov; Robert L. Devaney,
\emph{A First Course in Chaotic Dynamical Systems: Theory and Experiment},
2.\textsuperscript{a} ed., capítulos 4--8, sobre análisis gráfico, puntos fijos,
bifurcaciones y diagrama de órbitas; Kathleen T. Alligood, Tim D. Sauer y James A.
Yorke, \emph{Chaos: An Introduction to Dynamical Systems}, secciones 1.1--1.7.}

\subsection{Mapa de Hénon, plegamiento y dinámica simbólica}

\paragraph{Ecuaciones y puntos fijos de Hénon.}
Un mapa bidimensional puede contraer área en una dirección y estirar y plegar en
otra. El mapa de Hénon, en una convención habitual, es
\[
 x_{n+1}=1-a x_n^2+y_n,\qquad y_{n+1}=b x_n.
\]
Su Jacobiano es
\[
 DF(x,y)=
 \begin{pmatrix}-2ax&1\\ b&0\end{pmatrix},\qquad
 \det DF=-b.
\]
Por tanto, un elemento de área cambia su magnitud por el factor constante \(|b|\).
Con \(|b|<1\) existe contracción de área. La clasificación de cada órbita requiere
acotada.

La \term{dinámica simbólica} asigna a cada estado un símbolo según una partición. Un
itinerario \(s_0s_1s_2\cdots\) registra qué región visita la órbita. Para que esta
codificación sea una representación matemática completa se necesita justificar una
partición generadora y una relación con un desplazamiento simbólico; etiquetar puntos
por su signo sólo produce un resumen.

\paragraph{Ejemplo algebraico resuelto: puntos fijos de Hénon.}
En un punto fijo \(y^*=bx^*\) y
\[
 x^*=1-a(x^*)^2+bx^*,
\]
de modo que
\[
 a(x^*)^2+(1-b)x^*-1=0.
\]
Para \(a=1.4\), \(b=0.3\),
\[
 x_\pm^*=\frac{-0.7\pm\sqrt{6.09}}{2.8}
 \simeq 0.63135,\ -1.13135,
\qquad y_\pm^*=0.3x_\pm^*.
\]
En cada punto, \(\tau=-2ax^*\) y \(\Delta=-0.3<0\). Por el criterio plano, ambos
son sillas. El atractor numérico clásico, cuando se observa, pertenece a la dinámica alejada de esos puntos;
su organización está vinculada con las variedades de las sillas y el plegamiento
global.

\paragraph{Itinerario simbólico del mapa de Hénon.}
Defina \(s_n=L\) si \(x_n<0\) y \(s_n=R\) si \(x_n\geq0\). La palabra
\(RRL LR\), sin espacios, identifica cinco visitas, pero no conserva amplitudes ni
distancias. Las frecuencias de bloques \(LR\), \(RL\), \(LL\), \(RR\) pueden comparar
órbitas largas; la condición de bloque prohibido requiere muestras y criterios adicionales.

\begin{activitybox}{Práctica reproducible: Hénon y un alfabeto binario}
Genere Hénon con \(a=1.4\), \(b=0.3\), \((x_0,y_0)=(0,0)\), 20000 iteraciones y
descarte 2000. Use \gui{Series temporales} y \gui{Retratos 2D}; exporte los datos.
Verifique los dos puntos fijos calculados y clasifíquelos con sus matrices. Codifique
las últimas 10000 observaciones mediante \(L/R\), cuente bloques de longitud dos y
tres y repita con \((10^{-8},0)\). Registre si cambian las frecuencias dentro de la
precisión muestral.
\end{activitybox}

\begin{figure}[H]
\centering
\includegraphics[width=0.82\linewidth]{systems/henon_phase_timeseries.png}
\caption{Retrato y serie del mapa de Hénon generados por Toolbox Chaos. El retrato
revela plegamiento geométrico; la serie conserva el orden de visitas que se pierde en
la nube plana.}
\label{fig:henon-profundizacion}
\end{figure}

\paragraph{Lectura del retrato de fase de Hénon.}
El espesor aparente del conjunto depende de resolución, antialiasing y número de
puntos. Una ampliación debe recalcular o redibujar los datos, no ampliar una captura.
La serie permite detectar escape o periodicidad; el retrato revela foliaciones y
plegamiento; el itinerario simbólico resume recurrencia entre regiones.

\begin{warningbox}
\(|\det DF|<1\) demuestra contracción local de área. La existencia y unicidad de un
atractor, la sensibilidad y la conjugación con el desplazamiento de Bernoulli requieren verificar hipótesis adicionales sobre la partición y los itinerarios.
\end{warningbox}

\lecturaseccion{Steven H. Strogatz, \emph{Nonlinear Dynamics and Chaos with
Applications to Physics, Biology, Chemistry, and Engineering},
2.\textsuperscript{a} ed., sección 12.2, ``The Hénon Map''; Kathleen T. Alligood,
Tim D. Sauer y James A. Yorke, \emph{Chaos: An Introduction to Dynamical Systems},
secciones 2.2, 2.5--2.6, 3.5 y 5.6; Robert L. Devaney,
\emph{A First Course in Chaotic Dynamical Systems: Theory and Experiment},
2.\textsuperscript{a} ed., capítulo 9, sobre dinámica simbólica, desplazamiento y
conjugación; Morris W. Hirsch, Stephen Smale y Robert L. Devaney,
\emph{Differential Equations, Dynamical Systems, and an Introduction to Chaos},
3.\textsuperscript{a} ed., capítulo 15, apartados sobre caos, dinámica simbólica,
desplazamiento y conjunto de Cantor.}

\subsection{Lorenz, Rössler, Chua y Duffing--Ueda}

\paragraph{Mecanismos dinámicos en cuatro modelos.}
Estos sistemas son laboratorios matemáticos, no plantillas intercambiables. Lorenz
combina convección idealizada, simetría y disipación:
\[
 \dot x=\sigma(y-x),\qquad
 \dot y=x(\rho-z)-y,\qquad
 \dot z=xy-\beta z.
\]
Rössler concentra la no linealidad en una ecuación:
\[
 \dot x=-y-z,\qquad
 \dot y=x+ay,\qquad
 \dot z=b+z(x-c).
\]
Una forma de Chua usa una no linealidad por tramos:
\[
 \dot x=\alpha[y-x-h(x)],\quad
 \dot y=x-y+z,\quad
 \dot z=-\beta y,
\]
\[
 h(x)=m_1x+\frac{m_0-m_1}{2}
 \bigl(|x+1|-|x-1|\bigr).
\]
Duffing--Ueda es no autónomo; con la fase del forzamiento se escribe
\[
 \dot x=v,\quad
 \dot v=-2\zeta v-x-sx^3+f\cos\phi,\quad
 \dot\phi=\nu.
\]

\paragraph{Ejemplo algebraico resuelto: Lorenz clásico.}
La divergencia del campo es constante:
\[
 \nabla\cdot\mathbf f=-\sigma-1-\beta.
\]
Para \((\sigma,\rho,\beta)=(10,28,8/3)\) vale \(-41/3<0\); un volumen infinitesimal
se contrae como \(\exp[-(41/3)t]\). Esto demuestra disipación volumétrica local, no
acotación global por sí sola.

Los equilibrios son
\[
 E_0=(0,0,0),\qquad
 E_\pm=\bigl(\pm\sqrt{\beta(\rho-1)},
             \pm\sqrt{\beta(\rho-1)},\rho-1\bigr)
\]
cuando \(\rho>1\). En el caso clásico,
\[
 E_\pm=(\pm6\sqrt2,\pm6\sqrt2,27).
\]
La transformación \((x,y,z)\mapsto(-x,-y,z)\) deja invariantes las ecuaciones y
explica la simetría de los dos lóbulos. No obliga a una trayectoria finita a pasar
exactamente el mismo tiempo en cada uno.

\paragraph{Estructura algebraica de los cuatro modelos.}
En Rössler,
\[
 \nabla\cdot\mathbf f=a+x-c,
\]
que varía a lo largo de la órbita; en el Duffing autónomo,
\(\nabla\cdot\mathbf f=-2\zeta\). En Chua, la pendiente efectiva de \(h(x)\) cambia
al cruzar los quiebres. Por tanto, ``disipativo'' puede significar contracción
instantánea constante, contracción media o contracción por regiones; el cálculo debe
acompañar la etiqueta.

\begin{activitybox}{Práctica reproducible: atlas de cuatro flujos}
En cada pestaña seleccione el preset registrado de Lorenz, Rössler, Chua y
Duffing--Ueda. Antes de ejecutar copie las ecuaciones, parámetros y estado inicial que
muestra la interfaz. Use RK4 y el paso recomendado por cada preset; genere
\gui{Atractor 3D}, pulse \gui{Usar última trayectoria} en \gui{Retratos 2D} y
\gui{Series temporales}, y repita con la mitad del paso y el mismo horizonte físico.
Para Lorenz marque \(E_0,E_\pm\). Para Duffing interprete la fase como reloj del
forzamiento y construya, desde los datos, una sección cada \(2\pi/\nu\).
\end{activitybox}

\begin{figure}[H]
\centering
\includegraphics[width=0.90\linewidth]{systems/lorenz_phase_timeseries.png}
\caption{Lorenz en una vista combinada de trayectoria y series generada por Toolbox
Chaos. Los lóbulos reflejan la simetría del campo; los cambios de lóbulo se localizan
mejor al conservar el orden temporal.}
\label{fig:software-lorenz}
\end{figure}

\begin{figure}[H]
\centering
\includegraphics[width=0.90\linewidth]{systems/rossler_phase_timeseries.png}
\caption{Rössler generado por Toolbox Chaos. La rotación casi plana y las excursiones
en la tercera coordenada sugieren un mecanismo de estiramiento y reinyección; la
proyección requiere el procedimiento adicional que establece una sección de retorno unimodal.}
\label{fig:software-rossler}
\end{figure}

\begin{figure}[H]
\centering
\includegraphics[width=0.90\linewidth]{systems/chua_phase_timeseries.png}
\caption{Chua generado por Toolbox Chaos. Los cambios entre regiones de la
característica por tramos organizan la geometría de doble desplazamiento; los quiebres
de la ecuación deben interpretarse junto con la escala de los ejes.}
\label{fig:software-chua}
\end{figure}

\begin{figure}[H]
\centering
\includegraphics[width=0.90\linewidth]{systems/duffing_ueda_phase_timeseries.png}
\caption{Duffing--Ueda generado por Toolbox Chaos. El retrato proyectado mezcla
movimiento y fase de forzamiento; una sección estroboscópica separa mejor periodicidad,
cuasiperiodicidad y respuesta irregular.}
\label{fig:software-duffing}
\end{figure}

\paragraph{Variables para analizar cada modelo.}
En cada figura identifique: variable horizontal y vertical; tramo transitorio; rango
tardío; simetrías; vecindad de equilibrios; orientación; frecuencia de retornos; y
posibles saturaciones. Después formule una frase causal vinculada con un término de
la ecuación y una frase prudente que describa sólo lo observado. Esta doble redacción
evita que un parecido visual sustituya el análisis.

\begin{warningbox}
Las cuatro imágenes fueron generadas con configuraciones concretas del software. No
son retratos de todo el espacio de parámetros. La identificación de Lorenz, Rössler y caos en Chua requiere contrastar los retratos con ecuaciones y diagnósticos; una respuesta irregular de Duffing requiere controlar transitorios y aliasing.
\end{warningbox}

\lecturaseccion{Julien C. Sprott, \emph{Elegant Chaos: Algebraically Simple Chaotic
Flows}, capítulo 2, apartados sobre van der Pol, Rayleigh y Duffing forzados, y
capítulo 3, apartados sobre Lorenz, Rössler, sistemas de Sprott, ecuaciones jerk y
Chua; Kathleen T. Alligood, Tim D. Sauer y James A. Yorke,
\emph{Chaos: An Introduction to Dynamical Systems}, capítulo 9, apartados sobre
Lorenz, Rössler, Chua, flujos forzados y exponentes de Lyapunov; Morris W. Hirsch,
Stephen Smale y Robert L. Devaney, \emph{Differential Equations, Dynamical Systems,
and an Introduction to Chaos}, 3.\textsuperscript{a} ed., capítulos 14 y 16,
apartados sobre Lorenz, Rössler, doble desplazamiento y circuito de Chua.}

\subsection{Rutas al caos: duplicación, cuasiperiodicidad, intermitencia y crisis}

\paragraph{Mecanismos de transición al caos.}
Una \term{ruta al caos} es una secuencia organizada de cambios al variar un
parámetro. Entre las rutas frecuentes se encuentran:
\begin{itemize}
  \item duplicaciones sucesivas del periodo;
  \item ruptura de un toro después de movimiento cuasiperiódico;
  \item intermitencia, con fases casi regulares interrumpidas por ráfagas;
  \item crisis, donde una colisión global cambia bruscamente tamaño, acceso o vida de
  un conjunto caótico;
  \item bifurcaciones homoclínicas o heteroclínicas asociadas con variedades
  invariantes.
\end{itemize}
Estas etiquetas requieren evidencias diferentes. Un diagrama de máximos es eficaz
para duplicación y crisis; un número de rotación y espectro son más informativos para
cuasiperiodicidad; una distribución de longitudes laminares ayuda en intermitencia.

\paragraph{Ejemplo algebraico resuelto: primera cascada logística.}
La rama fija no nula del mapa logístico tiene multiplicador \(2-r\), que alcanza
\(-1\) en \(r=3\). Los puntos de periodo dos que no son fijos son
\[
 x_\pm=\frac{r+1\pm\sqrt{(r-3)(r+1)}}{2r}.
\]
Para \(r=3.2\),
\[
 x_-\simeq0.5130,\qquad x_+\simeq0.7995.
\]
El multiplicador del ciclo es el producto
\[
 M_2=f'(x_-)f'(x_+)=-r^2+2r+4.
\]
En \(r=3.2\), \(M_2=0.16\), por lo que el ciclo es estable. Pierde estabilidad
cuando \(M_2=-1\), es decir, en \(r=1+\sqrt6\simeq3.44949\). Este cálculo localiza
dos niveles de la cascada sin depender de la resolución del diagrama.

\paragraph{Escalamiento de Feigenbaum.}
Si \(r_n\) son parámetros de duplicaciones sucesivas en una familia unimodal
genérica, las razones
\[
 \delta_n=\frac{r_{n-1}-r_{n-2}}{r_n-r_{n-1}}
\]
pueden aproximarse a una constante universal. La universalidad requiere hipótesis
sobre la familia y un límite; una razón calculada con tres puntos aproximados no la
demuestra.

\begin{activitybox}{Práctica reproducible: distinguir cuatro rutas}
Primero use \gui{Bifurcación} para la logística y refine los intervalos alrededor de
\(r=3\) y \(r=1+\sqrt6\). Después, con Duffing--Ueda, construya secciones
estroboscópicas para una secuencia de amplitudes de forzamiento tomada del preset y
registre si aparece un número finito de puntos, una curva o una nube. Para Lorenz,
barra \(\rho\) en ambas direcciones y use el mismo estado inicial y transitorio.
Cuando observe un salto, simule a ambos lados con horizontes duplicados y varias
condiciones iniciales.
\end{activitybox}

\paragraph{Lectura de las rutas al caos.}
Una duplicación conserva orden: los máximos alternan entre dos, cuatro y más niveles.
Un toro se manifiesta como curva cerrada en una sección apropiada, no como una banda
en cualquier proyección. La intermitencia se reconoce en la serie por intervalos
laminares. Una crisis se investiga mediante cambio abrupto de tamaño o tiempo de vida
y proximidad a una estructura invariante.

\begin{warningbox}
Identifique una duplicación de periodo mediante la comprobación de la alternancia
temporal. Una banda espectral ensanchada no identifica por sí sola ruptura de toro.
Un salto en un barrido unidireccional puede ser retraso transitorio o coexistencia, y
un cambio abrupto requiere análisis adicional para clasificar una crisis global.
\end{warningbox}

\lecturaseccion{Steven H. Strogatz, \emph{Nonlinear Dynamics and Chaos with
Applications to Physics, Biology, Chemistry, and Engineering},
2.\textsuperscript{a} ed., secciones 8.4, 8.6--8.7 y 10.4--10.7; Robert C. Hilborn,
\emph{Chaos and Nonlinear Dynamics: An Introduction for Scientists and Engineers},
2.\textsuperscript{a} ed., capítulos 2, 4, 6 y 7, apartados sobre universalidad,
transitorios, cuasiperiodicidad, intermitencia y crisis; Kathleen T. Alligood, Tim D.
Sauer y James A. Yorke, \emph{Chaos: An Introduction to Dynamical Systems}, capítulos
10--12; John Guckenheimer y Philip Holmes, \emph{Nonlinear Oscillations, Dynamical
Systems, and Bifurcations of Vector Fields}, capítulos 4--6.}

\subsection{Fractales, cuencas y dimensiones}

\paragraph{Dimensiones fractales y cuencas.}
Un conjunto muestra \term{escalamiento fractal} cuando el número de elementos
necesarios para cubrirlo sigue una ley de potencia en un intervalo de escalas. La
dimensión de cajas se define, si existe el límite, por
\[
 D_B=\lim_{\varepsilon\to0}
 \frac{\log N(\varepsilon)}{\log(1/\varepsilon)}.
\]
Para puntos de una trayectoria, la integral de correlación
\[
 C(\varepsilon)=\frac{2}{N(N-1)}
 \sum_{1\leq i<j\leq N}
 \mathbf 1\{\|\mathbf x_i-\mathbf x_j\|<\varepsilon\}
\]
puede escalar como \(C(\varepsilon)\propto\varepsilon^{D_2}\). La estimación de
\(D_2\) requiere excluir pares temporalmente próximos, revisar saturación con la
dimensión de inmersión y encontrar una zona de escala.

\paragraph{Ejemplo algebraico resuelto: conjunto de Cantor.}
En la etapa \(k\) del conjunto ternario de Cantor hay \(N_k=2^k\) intervalos de
longitud \(\varepsilon_k=3^{-k}\). Por tanto,
\[
 D_B=\lim_{k\to\infty}\frac{\log2^k}{\log3^k}
 =\frac{\log2}{\log3}\simeq0.63093.
\]
La dimensión no proviene de que el dibujo sea rugoso, sino de una relación exacta de
autosemejanza.

\paragraph{Ejemplo espectral: Kaplan--Yorke.}
Ordenados \(\lambda_1\geq\cdots\geq\lambda_m\), sea \(j\) el mayor índice con
\(\sum_{i=1}^j\lambda_i\geq0\). La dimensión de Kaplan--Yorke se define como
\[
 D_{KY}=j+\frac{\sum_{i=1}^j\lambda_i}{|\lambda_{j+1}|}.
\]
Para un espectro ilustrativo \((0.90,0,-14.6)\),
\[
 D_{KY}=2+\frac{0.90}{14.6}\simeq2.0616.
\]
Ésta es una dimensión derivada de exponentes estimados, no la misma definición que
\(D_B\) o \(D_2\).

\begin{activitybox}{Práctica reproducible: cuenca y escalamiento}
En \gui{Cuenca de atracción}, seleccione un caso con clasificador disponible. Fije el
dominio de \((x_0,y_0)\), \(z_0\), integrador, \(dt\), horizonte y umbrales. Calcule
mallas \(100\times100\), \(200\times200\) y \(400\times400\); exporte la
clasificación. Para una frontera binaria, cuente cajas ocupadas a tamaños
\(\varepsilon=1/4,1/8,1/16,1/32\) del lado del dominio y ajuste la pendiente
\(\log N\) contra \(\log(1/\varepsilon)\). Repita con tiempo doble.
\end{activitybox}

\begin{figure}[H]
\centering
\includegraphics[width=0.84\linewidth]{lorenz_basin_reading_zones.png}
\caption{Zonas de lectura de una cuenca de Lorenz generada con Toolbox Chaos. Los
colores representan clases calculadas bajo un horizonte y un clasificador; las celdas
próximas a una frontera requieren refinamiento antes de asignarles etiquetas
asintóticas exactas.}
\label{fig:cuenca-profundizacion}
\end{figure}

\paragraph{Lectura de fronteras y dimensiones.}
Una pendiente sólo es defendible donde la gráfica log--log es aproximadamente lineal
y estable al añadir escalas. En una cuenca, las regiones interiores suelen ser más
robustas que la frontera. Conviene representar una clase separada para escapes,
indeterminación o no convergencia, en vez de forzar cada celda a un atractor.

\begin{warningbox}
Una dimensión estimada con dos escalas es sólo una pendiente. El antialiasing de una
imagen no contiene datos de cajas; deben usarse coordenadas o etiquetas exportadas.
La identificación de Wada requiere pruebas topológicas, y un valor de Kaplan--Yorke se complementa con una
estimación geométrica y evidencia de la medida invariante específica.
\end{warningbox}

\lecturaseccion{Kathleen T. Alligood, Tim D. Sauer y James A. Yorke,
\emph{Chaos: An Introduction to Dynamical Systems}, secciones 4.1, 4.4--4.7 y 5.3;
Robert L. Devaney, \emph{A First Course in Chaotic Dynamical Systems: Theory and
Experiment}, 2.\textsuperscript{a} ed., capítulo 14, apartados sobre fractales y
dimensión; Robert C. Hilborn, \emph{Chaos and Nonlinear Dynamics: An Introduction
for Scientists and Engineers}, 2.\textsuperscript{a} ed., capítulos 9--10, apartados
sobre dimensión fractal, correlación y multifractales; Tamás Tél y Márton Gruiz,
\emph{Chaotic Dynamics: An Introduction Based on Classical Mechanics}, capítulos
2 y 6.}

\subsection{Series temporales, FFT y densidad espectral}

\paragraph{Series temporales, FFT y PSD.}
Una serie uniformemente muestreada es \(s_n=s(n\Delta t)\), con frecuencia
\(f_s=1/\Delta t\). La transformada discreta
\[
 S_k=\sum_{n=0}^{N-1}s_n e^{-2\pi i kn/N}
\]
evalúa frecuencias \(f_k=kf_s/N\). Para una señal real, la parte unilateral contiene
frecuencias hasta Nyquist,
\[
 f_N=\frac{f_s}{2}.
\]
El \term{espectro de amplitud} describe magnitudes de componentes en un registro; la
\term{densidad espectral de potencia} distribuye potencia por frecuencia. Welch
promedia periodogramas de segmentos solapados y reduce varianza a cambio de resolución.

\paragraph{Ventanas, fuga y aliasing.}
Un registro finito equivale a multiplicar la señal por una ventana temporal. Si una
sinusoide no completa un número entero de periodos, su energía se reparte entre
frecuencias: \term{fuga espectral}. Una ventana suaviza discontinuidades en los
extremos, pero ensancha picos. Si existe contenido por encima de \(f_N\), el muestreo
lo refleja a frecuencias falsas: \term{aliasing}. Ninguna opción gráfica repara datos
ya submuestreados.

\paragraph{Ejemplo de frecuencias discretas.}
Con \(\Delta t=0.01\) y \(N=10000\),
\[
 f_s=100,\qquad f_N=50,\qquad
 \Delta f=\frac{f_s}{N}=0.01.
\]
Una señal \(s(t)=2\cos(2\pi\,1.5\,t)\) cae exactamente en el bin
\(k=1.5/0.01=150\) si el registro mantiene esa longitud. Si se reduce a
\(N=2000\), \(\Delta f=0.05\); el pico sigue localizable pero con malla cinco veces
más gruesa. Una señal de \(60\) Hz muestreada a \(100\) Hz aparece aliased cerca de
\(|60-100|=40\) Hz.

\begin{activitybox}{Práctica reproducible: una trayectoria, dos espectros}
Genere Lorenz clásico con RK4, \(dt=0.01\), \(T=200\), y descarte \(T_0=20\).
Pulse \gui{Usar última trayectoria} en \gui{Espectro}. Calcule primero
\gui{PSD de Welch (recomendado)} y después \gui{Espectro de amplitud}, sin cambiar
datos ni variable. Repita con \(T_0=50\) y con una trayectoria generada a
\(dt=0.005\). Exporte las series y compruebe \(f_s\), Nyquist, longitud útil y
resolución. Compare \(x\), \(y\) y \(z\) junto con el pico más alto.
\end{activitybox}

\begin{figure}[H]
\centering
\includegraphics[width=0.80\linewidth]{lorenz_fft.png}
\caption{Espectro normalizado centrado de Lorenz generado a partir de datos del software. Los
picos y la banda describen el registro finito y su muestreo; su interpretación requiere
solos una prueba de caos.}
\label{fig:fft-profundizacion}
\end{figure}

\paragraph{Lectura de espectros y series.}
Una órbita periódica ideal produce líneas en la frecuencia fundamental y armónicos;
una cuasiperiódica, combinaciones de varias frecuencias inconmensurables; una señal
irregular puede mostrar banda ancha. Transitorios, ruido, modulación y no
estacionariedad también ensanchan el espectro. La serie original debe revisarse antes
de interpretar cada pico.

\begin{warningbox}
Un espectro ancho orienta el análisis de caos y un pico agudo orienta el de periodicidad
asintótica; ambos requieren pruebas dinámicas complementarias. Comparar amplitudes entre registros con distinta normalización, ventana o
duración es inválido. Welch mejora estabilidad estadística, pero no aumenta la
frecuencia de Nyquist ni recupera información perdida por aliasing.
\end{warningbox}

\lecturaseccion{Robert C. Hilborn, \emph{Chaos and Nonlinear Dynamics: An
Introduction for Scientists and Engineers}, 2.\textsuperscript{a} ed., apéndice A,
``Fourier Analysis'', y capítulo 9, apartados sobre cuantificación de series; Julien
C. Sprott, \emph{Elegant Chaos: Algebraically Simple Chaotic Flows}, capítulo 1,
apartados sobre espectros, ciclos, quasiperiodicidad e indicadores; Kathleen T.
Alligood, Tim D. Sauer y James A. Yorke, \emph{Chaos: An Introduction to Dynamical
Systems}, capítulo 13, apartado inicial sobre datos de series temporales.}

\subsection{Reconstrucción por retardo y diagnóstico desde una observación}

\paragraph{Vectores de reconstrucción por retardo.}
Cuando sólo se mide un escalar \(s(t)\), se forman vectores de retardo
\[
 \mathbf y_n=
 \bigl(s_n,s_{n-\ell},s_{n-2\ell},\ldots,s_{n-(m-1)\ell}\bigr),
\]
donde \(\tau=\ell\Delta t\) es el retardo y \(m\) la dimensión de inmersión. Bajo
hipótesis de suavidad, genericidad y dimensión suficiente, un teorema de inmersión
permite reconstruir una geometría equivalente en un sentido preciso. En datos
finitos, elegir \(\tau\), \(m\), observable y preprocesamiento es parte del resultado.

\paragraph{Selección de retardo y dimensión.}
Un retardo demasiado pequeño agrupa puntos cerca de la diagonal; uno demasiado grande
descorrelaciona las componentes y mezcla escalas. Se usan, como guías, el primer cruce
o mínimo de autocorrelación o información mutua. La dimensión puede explorarse con
falsos vecinos cercanos o saturación de una magnitud. Ningún umbral es universal.

\paragraph{Ejemplo algebraico resuelto: oscilador armónico.}
Sea \(s(t)=\cos(\omega t)\), con periodo \(T=2\pi/\omega\). Si
\(\tau=T/4\),
\[
 s(t-\tau)=\cos\left(\omega t-\frac{\pi}{2}\right)=\sin(\omega t),
\]
y por tanto
\[
 [s(t)]^2+[s(t-\tau)]^2=1.
\]
La reconstrucción bidimensional es un círculo. Si \(\tau=0\), todos los puntos caen
en la diagonal; si \(\tau=T/2\), caen en la antidiagonal. El mismo sistema puede
parecer degenerado por una elección inadecuada del retardo.

\begin{activitybox}{Práctica reproducible: reconstruir Lorenz desde \(x(t)\)}
En Toolbox Chaos genere Lorenz con \(dt=0.01\), \(T=200\) y descarte \(20\).
Exporte \(x(t)\). Calcule la autocorrelación y elija tres retardos: uno corto, el
primer cruce o escala equivalente, y uno cuatro veces mayor. Construya las nubes
\((x_n,x_{n-\ell})\) y
\((x_n,x_{n-\ell},x_{n-2\ell})\). Registre número de muestras efectivas después de
los desplazamientos. Repita con \(z(t)\) y explique qué rasgos dependen del observable.
\end{activitybox}

\paragraph{Lectura de reconstrucciones por retardo.}
Una reconstrucción útil separa estados sin convertir el conjunto en nube uniforme.
Las hebras cercanas pueden ser estructura o escasez de datos; deben revisarse con otro
retardo y más muestras. Un retrato de retorno \(s_{n+1}\) contra \(s_n\) aporta una representación distinta de
una inmersión general: sólo usa retraso de una muestra.

\begin{warningbox}
Una nube con forma de atractor no verifica las hipótesis de un teorema de inmersión.
La estimación de dimensión requiere puntos temporalmente separados, y el determinismo requiere más evidencia que una reconstrucción suave. Filtrar, normalizar o
interpolar cambia distancias y debe documentarse.
\end{warningbox}

\lecturaseccion{Kathleen T. Alligood, Tim D. Sauer y James A. Yorke,
\emph{Chaos: An Introduction to Dynamical Systems}, capítulo 13, apartados sobre
datos, coordenadas de retardo y embedología; Robert C. Hilborn,
\emph{Chaos and Nonlinear Dynamics: An Introduction for Scientists and Engineers},
2.\textsuperscript{a} ed., capítulo 10, apartados sobre inmersión, reconstrucción y
multifractales; Steven H. Strogatz, \emph{Nonlinear Dynamics and Chaos with
Applications to Physics, Biology, Chemistry, and Engineering},
2.\textsuperscript{a} ed., sección 12.4, ``Reconstructing Attractors from Data''.}

\subsection{Sistemas con retardo, alta dimensión e hipercaos}

\paragraph{Historia de estado en Mackey--Glass.}
En una ecuación con retardo,
\[
 \dot x(t)=F\bigl(x(t),x(t-\tau)\bigr),
\]
el valor \(x(0)\) se complementa con una función de historia que determina el futuro. Debe especificarse una función
historia \(\varphi(t)\) en \([-\tau,0]\). El espacio de estados es entonces
infinito-dimensional, aunque la salida se represente con unas pocas coordenadas de
retardo.

\paragraph{Ejemplo algebraico resuelto: Mackey--Glass.}
Una forma habitual es
\[
 \dot x(t)=
 \frac{\beta x(t-\tau)}{1+x^n(t-\tau)}-\gamma x(t).
\]
Un equilibrio constante \(x(t)=x^*\) satisface
\[
 0=x^*\left[\frac{\beta}{1+(x^*)^n}-\gamma\right].
\]
Además de \(x^*=0\), existe
\[
 x^*=\left(\frac{\beta}{\gamma}-1\right)^{1/n}
\]
si \(\beta>\gamma\). Para \(\beta=0.2\), \(\gamma=0.1\), \(n=10\),
\(x^*=1\). El equilibrio no depende de \(\tau\), pero su estabilidad sí: el retardo
entra en la ecuación característica trascendente de la linealización.

\paragraph{Alta dimensión: Lorenz--96.}
Con índices cíclicos,
\[
 \dot x_i=(x_{i+1}-x_{i-2})x_{i-1}-x_i+F,
 \qquad i=1,\ldots,N.
\]
El estado homogéneo \(x_i=F\) es equilibrio porque el término de advección se anula.
La divergencia es
\[
 \sum_{i=1}^{N}\frac{\partial\dot x_i}{\partial x_i}=-N,
\]
de modo que el modelo contrae volumen en el espacio \(N\)-dimensional. Alta dimensión
requiere criterios adicionales para \term{hipercaos}; esta etiqueta suele requerir más de un exponente de
Lyapunov positivo convergente.

\begin{activitybox}{Práctica reproducible: historia y proyección}
Ejecute el preset de Mackey--Glass y registre \(\beta,\gamma,n,\tau\), historia,
paso y esquema numérico mostrado por la interfaz. Duplique el horizonte y compare
series y retratos de retardo. Después ejecute Lorenz--96 con la dimensión y forzamiento
del preset; grafique tres componentes contiguas y tres separadas. Exporte las series y
compare correlaciones y escalas. Si \gui{Lyapunov} aparece como no disponible,
consérvelo como límite del contrato y no transfiera el algoritmo 3D.
\end{activitybox}

\begin{figure}[H]
\centering
\includegraphics[width=0.88\linewidth]{systems/mackey_glass_phase_timeseries.png}
\caption{Mackey--Glass generado por Toolbox Chaos. La proyección usa retardos o
componentes derivadas de una historia; conserva la naturaleza funcional de la ecuación y requiere el marco adecuado de EDO
tridimensional.}
\label{fig:mackey-glass-profundizacion}
\end{figure}

\begin{figure}[H]
\centering
\includegraphics[width=0.88\linewidth]{systems/lorenz96_phase_timeseries.png}
\caption{Lorenz--96 generado por Toolbox Chaos. Tres coordenadas muestran una sombra
del estado de dimensión mayor; la estructura omitida debe discutirse antes de inferir
geometría global.}
\label{fig:lorenz96-profundizacion}
\end{figure}

\paragraph{Lectura de sistemas con retardo y alta dimensión.}
En retardo, una transición depende de la historia junto con el valor final.
En alta dimensión, dos proyecciones distintas pueden ocultar o crear cruces aparentes.
La presencia de varias escalas temporales sugiere calcular correlaciones, espectros y,
cuando exista un método compatible, varios exponentes.

\begin{warningbox}
Use una historia completa para representar el estado. Clasifique como hipercaótico
un sistema por tener cuatro o más variables, ni extrapole un espectro de Lyapunov 3D
a dimensión alta. Una señal irregular de Mackey--Glass o Lorenz--96 puede incluir un
transitorio largo y requiere prueba de horizonte.
\end{warningbox}

\lecturaseccion{Julien C. Sprott, \emph{Elegant Chaos: Algebraically Simple Chaotic
Flows}, capítulo 6, apartados sobre sistemas de alta dimensión e hipercaos, capítulo
8, ``Spatiotemporal Chaos'', y capítulo 9, ``Delay Differential Equations''; Robert
C. Hilborn, \emph{Chaos and Nonlinear Dynamics: An Introduction for Scientists and
Engineers}, 2.\textsuperscript{a} ed., capítulo 11, apartados sobre caos
espaciotemporal; Tamás Tél y Márton Gruiz, \emph{Chaotic Dynamics: An Introduction
Based on Classical Mechanics}, capítulo 3, apartados sobre estabilidad y espacio de
fases general.}

\subsection{Caos conservativo y caos transitorio}

\paragraph{Estructura de los sistemas conservativos.}
En un sistema hamiltoniano
\[
 \dot q=\frac{\partial H}{\partial p},\qquad
 \dot p=-\frac{\partial H}{\partial q},
\]
se cumple \(\dot H=0\) y la divergencia del campo es cero. No hay atracción
volumétrica como en Lorenz. El espacio de fases puede mezclar islas regulares, toros y
mares caóticos; hablar de atractor en el sentido disipativo suele ser inapropiado.

\paragraph{Ejemplo algebraico resuelto: mapa estándar.}
El mapa
\[
 p_{n+1}=p_n+K\sin q_n,\qquad
 q_{n+1}=q_n+p_{n+1}\pmod{2\pi}
\]
tiene Jacobiano
\[
 D F=
 \begin{pmatrix}
 1+K\cos q_n&1\\
 K\cos q_n&1
 \end{pmatrix},
\qquad \det DF=1.
\]
Por tanto preserva área exactamente en aritmética matemática. Para \(K=0\),
\(p_n=p_0\) y \(q_n=q_0+np_0\); el número de rotación es \(p_0/(2\pi)\). Al aumentar
\(K\), resonancias y capas caóticas pueden coexistir sin contracción global.

\paragraph{Tiempo de escape y caos transitorio.}
Una silla caótica no atractora puede retener órbitas durante tiempos irregulares antes
de que escapen. Si la probabilidad de supervivencia se aproxima por
\[
 P(t)\simeq e^{-\kappa t},
\]
\(\kappa\) es una tasa de escape y el tiempo medio idealizado es \(1/\kappa\). Si de
1000 semillas quedan cerca de 368 a \(t=10\), entonces
\[
 \widehat\kappa=-\frac1{10}\log(368/1000)\simeq0.10,
\qquad \widehat{\langle T\rangle}\simeq10.
\]
La estimación de esta distribución requiere un conjunto de trayectorias y condiciones iniciales.

\begin{activitybox}{Práctica reproducible: conservación y supervivencia}
En \gui{Crear sistema}, defina el oscilador
\(\dot q=p,\ \dot p=-q\) con \((q_0,p_0)=(1,0)\). Para Euler, Heun y RK4 y
\(h=0.05,0.025,0.0125\), exporte \(H_n=(q_n^2+p_n^2)/2\) y grafique su deriva.
Después seleccione un caso del catálogo que muestre escape o transitorio prolongado,
si está disponible, y ejecute una malla explícita de condiciones iniciales con dos
horizontes. Construya una curva de supervivencia sin reclasificar escapes como
atractores.
\end{activitybox}

\paragraph{Lectura de trayectorias conservativas y transitorias.}
En dinámica conservativa, puntos que llenan una curva indican movimiento regular
confinado; nubes e islas pueden coexistir. La deriva sistemática de energía revela
efecto del integrador, no disipación física. En caos transitorio, se debe graficar
número de sobrevivientes en escala semilogarítmica y comprobar si existe un intervalo
aproximadamente lineal.

\begin{warningbox}
En un sistema conservativo, una trayectoria acotada debe analizarse mediante invariantes y estructura de volumen antes de clasificarla como atractor. Un determinante
numérico cercano a uno se complementa con controles de conservación. El horizonte finito
censura tiempos de escape, por lo que las irregularidades requieren extender la simulación y comparar condiciones iniciales.
\end{warningbox}

\lecturaseccion{Robert C. Hilborn, \emph{Chaos and Nonlinear Dynamics: An
Introduction for Scientists and Engineers}, 2.\textsuperscript{a} ed., capítulo 4,
apartados sobre transitorios caóticos y órbitas homoclínicas, y capítulo 8, apartados
sobre dinámica hamiltoniana, KAM, mapa estándar y Hénon--Heiles; Tamás Tél y Márton
Gruiz, \emph{Chaotic Dynamics: An Introduction Based on Classical Mechanics},
capítulo 6, ``Transient Chaos'', y capítulo 7, apartados sobre dinámica conservativa,
mapa estándar, ergodicidad y mezcla; Julien C. Sprott, \emph{Elegant Chaos:
Algebraically Simple Chaotic Flows}, capítulo 4, apartados sobre flujos conservativos.}

\subsection{Sincronización y control de caos}

\paragraph{Sincronización y variables de error.}
Dos sistemas se \term{sincronizan} cuando una medida de su diferencia converge o
permanece pequeña bajo un acoplamiento especificado. La sincronía completa usa
\(\|\mathbf x_1-\mathbf x_2\|\to0\); existen también sincronía de fase, generalizada
y retrasada. \term{Control de caos} introduce una intervención pequeña y diseñada
para estabilizar una órbita o región deseada. La evaluación del control declara señal, ley de control, costo y estabilidad del lazo cerrado.

\paragraph{Dinámica transversal del error.}
Para dos sistemas idénticos con acoplamiento difusivo completo,
\[
 \dot{\mathbf x}_1=\mathbf f(\mathbf x_1),\qquad
 \dot{\mathbf x}_2=\mathbf f(\mathbf x_2)+k(\mathbf x_1-\mathbf x_2),
\]
el error \(\mathbf e=\mathbf x_2-\mathbf x_1\) satisface, cerca de la variedad
sincrónica,
\[
 \dot{\mathbf e}\simeq
 [D\mathbf f(\mathbf x_1(t))-kI]\mathbf e.
\]
La estabilidad depende de la expansión transversal a lo largo de la trayectoria, no
sólo de los autovalores en un equilibrio.

\paragraph{Ejemplo algebraico resuelto.}
Para el par lineal
\[
 \dot x=-ax,\qquad \dot y=-ay+k(x-y),
\]
\(e=y-x\) cumple
\[
 \dot e=-(a+k)e,\qquad e(t)=e(0)e^{-(a+k)t}.
\]
Si \(a=0.2\), \(k=0.8\), el error disminuye con tasa \(1\). Sin acoplamiento,
\(k=0\), disminuye con tasa \(0.2\). Este ejemplo separa sincronización inducida de
la convergencia común a un equilibrio; en un sistema caótico, cada trayectoria puede
seguir siendo irregular mientras el error transversal decrece.

\paragraph{Control en un mapa linealizado.}
Si cerca de una órbita objetivo el mapa controlado es
\[
 \xi_{n+1}=a\xi_n+b u_n,\qquad u_n=-K\xi_n,
\]
el multiplicador cerrado es \(a-bK\). Para \(a=1.2\), \(b=1\), elegir \(K=0.5\)
produce \(0.7\), estable. Debe comprobarse que la linealización y los límites del
actuador sean válidos en la región usada.

\begin{activitybox}{Práctica reproducible: error de dos Lorenz}
En \gui{Crear sistema}, defina dos copias de Lorenz con variables
\texttt{x1,y1,z1,x2,y2,z2}. La primera usa las ecuaciones clásicas; en la segunda
añada \texttt{k*(x1-x2)}, \texttt{k*(y1-y2)} y \texttt{k*(z1-z2)} a las tres
ecuaciones. Use \((10,28,8/3)\), estados \((1,1,1)\) y \((1.001,1,1)\),
\(k=0,1,5\), RK4, \(h=0.005\), \(T=50\). Exporte datos y grafique
\(\log\|\mathbf e(t)\|\), evitando el logaritmo de cero. Repita con \(h/2\) y
registre la magnitud del acoplamiento.
\end{activitybox}

\paragraph{Lectura del error de sincronización.}
La sincronización se evalúa en el error, no por superposición visual de dos curvas
gruesas. Una recta descendente en escala semilogarítmica sugiere contracción
exponencial transversal en ese intervalo. En control, compare régimen sin control,
transitorio de activación, estado controlado y esfuerzo \(u(t)\).

\begin{warningbox}
La estabilidad global de la variedad sincrónica requiere controles más amplios que un
error pequeño durante tiempo finito. Las pruebas de robustez usan perturbaciones de las condiciones iniciales. Un controlador que satura, cambia mucho el modelo o
requiere conocer toda la trayectoria debe reportarse con esas condiciones de uso;
estabilización de pequeña señal.
\end{warningbox}

\lecturaseccion{Robert C. Hilborn, \emph{Chaos and Nonlinear Dynamics: An
Introduction for Scientists and Engineers}, 2.\textsuperscript{a} ed., capítulo 12,
apartados sobre control y sincronización; Julien C. Sprott,
\emph{Elegant Chaos: Algebraically Simple Chaotic Flows}, capítulo 1, apartados sobre
exponentes, estabilidad y transitorios, y capítulo 10, apartados sobre realizaciones
en circuitos; John Guckenheimer y Philip Holmes, \emph{Nonlinear Oscillations,
Dynamical Systems, and Bifurcations of Vector Fields}, capítulo 4, apartados sobre
perturbaciones y métodos de Melnikov, como contexto de intervención débil.}

\section{Proyecto final sobre un sistema dinámico}

\paragraph{Pregunta y evidencia del proyecto.}
El proyecto integra un resultado analítico con evidencia numérica. La pregunta debe
ser acotada, por ejemplo: ¿cómo cambia un régimen al variar un parámetro?, ¿qué
condiciones iniciales alcanzan dos destinos?, ¿qué propiedad permanece al refinar el
paso?, ¿qué observable distingue periodicidad de irregularidad? El informe acepta como
pregunta únicamente ``graficar un atractor''.

\paragraph{Opciones de proyecto.}
\begin{enumerate}
  \item \textbf{Mapa:} logística, Hénon o sistema del \gui{Explorador Sprott};
  incluir puntos periódicos, multiplicadores, bifurcación y dinámica simbólica.
  \item \textbf{Flujo clásico:} Lorenz, Rössler, Chua o Duffing--Ueda; incluir
  equilibrios, Jacobiano, disipación, sección y al menos un indicador.
  \item \textbf{Cuencas y coexistencia:} comparar condiciones iniciales, resolución,
  horizonte y clasificador, con categoría indeterminada.
  \item \textbf{Datos:} reconstruir una dinámica desde una serie, justificar muestreo,
  FFT/PSD, retardo, dimensión y control sustituto.
  \item \textbf{Extensión:} retardo, alta dimensión, caos transitorio,
  sincronización o memoria fraccionaria, con alcance explícitamente prudente.
\end{enumerate}

\paragraph{Desarrollo algebraico obligatorio.}
El informe debe incluir al menos uno de los siguientes cálculos completos:
\begin{itemize}
  \item puntos fijos o equilibrios y clasificación mediante derivada o Jacobiano;
  \item adimensionalización y grupos de control;
  \item multiplicador de órbita o forma normal local;
  \item divergencia, invariante o balance de energía;
  \item ecuación variacional y significado del indicador de Lyapunov;
  \item solución de equilibrio para un sistema con retardo;
  \item dinámica del error para sincronización o control.
\end{itemize}
Copiar ecuaciones del catálogo sin desarrollar una consecuencia no satisface este
requisito.

\paragraph{Protocolo computacional mínimo.}
\begin{enumerate}
  \item Registrar nombre del sistema, ecuaciones, parámetros y unidades o escalas.
  \item Guardar condición o historia inicial, método, paso, tolerancias, horizonte,
  transitorio, muestreo y umbral de escape.
  \item Producir tres vistas coordinadas desde los mismos datos cuando sea posible.
  \item Exportar datos y metadatos junto con las imágenes.
  \item Repetir con al menos dos pasos u opciones de tolerancia.
  \item Repetir con horizonte o transitorio distinto.
  \item Comparar condiciones iniciales si la pregunta es global.
  \item Explicar disponibilidad y limitación de cada pestaña utilizada.
\end{enumerate}

\begin{evidencebox}{Matriz de afirmaciones del proyecto}
Para cada conclusión, incluya cuatro columnas: afirmación; evidencia directa;
control de robustez; y alternativa no descartada. Ejemplo: ``la estimación finita de
\(\lambda_1\) es positiva'' se apoya en la curva de convergencia y dos refinamientos;
describe el conjunto de estados y parámetros que cubre la evidencia declarada.
\end{evidencebox}

\paragraph{Estructura del informe.}
\begin{longtable}{>{\raggedright\arraybackslash}p{0.23\linewidth}
                  >{\raggedright\arraybackslash}p{0.67\linewidth}}
\toprule
Sección & Contenido verificable \\
\midrule
\endhead
Pregunta y alcance & Variable de respuesta, parámetro o comparación y límites.\\
Modelo & Ecuaciones, estado, hipótesis, unidades y procedencia del preset.\\
Análisis & Desarrollo algebraico y predicción anterior a la simulación.\\
Método & Contrato numérico completo y versión del software.\\
Resultados & Datos, tablas, tres vistas, pies autosuficientes e indicador.\\
Robustez & Paso, horizonte, transitorio, semilla y, cuando proceda, resolución.\\
Discusión & Coincidencias y tensiones entre predicción, gráfica e indicador.\\
Conclusión & Respuesta graduada y lista de propiedades no demostradas.\\
Anexo & Archivos JSON/CSV, nombres de figuras y guía de reproducción.\\
\bottomrule
\end{longtable}

\begin{activitybox}{Reproducción cruzada}
Entregue a otro equipo una carpeta que contenga únicamente informe, configuración y
datos. El equipo receptor repite una figura y un valor cuantitativo sin explicación
oral. Debe registrar versión, diferencias y ambigüedades. El equipo autor corrige el
protocolo y conserva ambas ejecuciones como evidencia de trazabilidad.
\end{activitybox}

\begin{warningbox}
El proyecto formula conclusiones acordes con muestreos finitos. Las expresiones ``demuestra caos'', ``prueba estabilidad global'' o ``valida el
modelo físico'' requieren un argumento matemático específico con todas sus hipótesis y evidencia adicional a una pestaña o una figura.
\end{warningbox}

\lecturaseccion{Edward N. Lorenz, \emph{The Essence of Chaos}, capítulos 1--5 y
apéndices ``The Butterfly Effect'', ``Mathematical Excursions'' y ``Glossary'', como
modelo de conexión entre mecanismo, cálculo e interpretación; Robert C. Hilborn,
\emph{Chaos and Nonlinear Dynamics: An Introduction for Scientists and Engineers},
2.\textsuperscript{a} ed., capítulos 9--10, sobre cuantificación y reconstrucción;
Kathleen T. Alligood, Tim D. Sauer y James A. Yorke,
\emph{Chaos: An Introduction to Dynamical Systems}, apéndice B, para control de error
numérico.}

\section[\texorpdfstring{\hspace{0.45em}}{}Banco graduado de actividades y problemas]{Banco graduado de actividades y problemas}

\paragraph{Estructura de actividades A--C.}
Las actividades A son de interpretación, las B combinan álgebra y simulación, y las C
constituyen miniproyectos. En todos los casos se entrega predicción previa, contrato
numérico, figura con ejes, datos exportados, control y evidencia pendiente por documentar.

\subsection*{Nivel A: lectura y conceptos}
\begin{enumerate}
  \item[A1.] \textbf{Misma órbita, tres vistas.} Use Lorenz y marque en
  \gui{Atractor 3D}, \gui{Retratos 2D} y \gui{Series temporales} el mismo intervalo.
  Explique qué información pierde cada vista.
  \item[A2.] \textbf{Mapa y flujo: reglas de actualización.} Para veinte puntos logísticos, compare una
  gráfica con marcadores y otra con líneas. Identifique qué segmentos no pertenecen a
  la dinámica.
  \item[A3.] \textbf{Nullcline y trayectoria: funciones complementarias.} Sobre el ejemplo
  \(\dot x=x(1-x-y)\), \(\dot y=y(x-1/2)\), señale dónde el vector es vertical y
  dónde horizontal.
  \item[A4.] \textbf{Transitorio.} Para Rössler, grafique las primeras y las últimas
  veinte unidades de tiempo con la misma escala. Decida qué estadísticas cambian.
  \item[A5.] \textbf{Proyección engañosa.} Localice un cruce en \(x\)-\(y\) de
  Lorenz y compare las coordenadas \(z\) y los tiempos de ambos puntos.
  \item[A6.] \textbf{Pie de figura auditable.} Reescriba el pie de una captura para
  incluir sistema, parámetros, estado, método, paso, tiempo, transitorio y mensaje
  principal.
\end{enumerate}

\subsection*{Nivel B: cálculo y reproducción}
\begin{enumerate}
  \item[B1.] \textbf{Saddle--node.} Derive ramas y estabilidad de
  \(\dot x=\mu-x^2\); mida el tiempo para entrar en una tolerancia de \(10^{-3}\)
  cuando \(\mu\) se aproxima a cero.
  \item[B2.] \textbf{Plano traza--determinante.} Diseñe una matriz para cada región,
  calcule sus autovalores y reproduzca los retratos en \gui{Crear sistema}.
  \item[B3.] \textbf{Jacobiano de Lorenz.} Derive \(J(x,y,z)\), evalúelo en
  \(E_0,E_\pm\) y confronte \gui{Autovalores} para dos valores de \(\rho\).
  \item[B4.] \textbf{Ciclo radial.} Estime periodo y multiplicador del ciclo
  \(r=1\) con dos pasos; compare con \(2\pi\) y \(e^{-4\pi}\).
  \item[B5.] \textbf{Periodo dos logístico.} Para \(r=3.2\), verifique
  \(x_\pm\), su orden temporal y el multiplicador \(M_2=0.16\).
  \item[B6.] \textbf{Contracción de Hénon.} Seleccione tres triángulos pequeños de
  condiciones iniciales, itere una vez y compare áreas con el factor \(|b|\).
  \item[B7.] \textbf{Divergencia.} Calcule \(\nabla\cdot f\) para Lorenz, Rössler y
  Duffing autónomo; relacione signo instantáneo con el volumen, sin inferir cuenca.
  \item[B8.] \textbf{Convergencia Lyapunov.} Para Lorenz, combine dos pasos, dos
  tiempos finales y dos intervalos QR. Informe sólo dígitos estables.
  \item[B9.] \textbf{Nyquist.} Genere dos señales sintéticas en
  \gui{Crear sistema} o desde datos controlados, una debajo y otra encima de Nyquist;
  identifique el alias en \gui{Espectro}.
  \item[B10.] \textbf{Reconstrucción.} Para una señal periódica, compruebe que
  \(\tau=T/4\) produce un círculo y que \(\tau=0\) produce una diagonal.
  \item[B11.] \textbf{Dimensión de cajas.} Valide el algoritmo con etapas del Cantor
  antes de aplicarlo a una frontera de cuenca.
  \item[B12.] \textbf{Energía numérica.} Compare deriva de \(H\) para Euler, Heun y
  RK4; interprete la menor deriva como un indicador que se complementa con la prueba de conservación.
\end{enumerate}

\subsection*{Nivel C: miniproyectos}
\begin{enumerate}
  \item[C1.] \textbf{Ventana periódica.} Localice una ventana logística, verifique
  periodo, multiplicador y \(\lambda_N\), y estudie sus bordes.
  \item[C2.] \textbf{Retorno de Rössler.} Defina una sección, interpole cruces,
  construya \(u_{n+1}\) contra \(u_n\) y discuta si es aproximadamente unimodal.
  \item[C3.] \textbf{Ruta de Duffing.} Varíe un solo parámetro, compare secciones
  estroboscópicas y clasifique la evidencia de periodicidad, toro o irregularidad.
  \item[C4.] \textbf{Frontera de cuenca.} Refine resolución y horizonte, estime una
  pendiente y cuantifique celdas que cambian de clase.
  \item[C5.] \textbf{Mackey--Glass.} Compare dos historias con el mismo \(x(0)\);
  muestre cómo la historia completa define el estado.
  \item[C6.] \textbf{Lorenz--96.} Compare proyecciones contiguas y separadas, PSD y
  correlaciones; explique qué queda fuera de la vista.
  \item[C7.] \textbf{Caos transitorio.} Construya una distribución de tiempos de
  escape con datos censurados y pruebe un ajuste exponencial sólo en el intervalo
  sustentado.
  \item[C8.] \textbf{Sincronización.} Estime una tasa transversal desde
  \(\log\|\mathbf e(t)\|\), explore \(k\) y cuantifique el costo de acoplamiento.
  \item[C9.] \textbf{Acceso a coexistencia.} Compare semillas cercanas a equilibrios
  y alejadas; clasifique la evidencia sin convertir ausencia finita en ocultamiento.
  \item[C10.] \textbf{Explorador Sprott.} Elija un código sintético público,
  decodifique sus coeficientes, derive sus ecuaciones, calcule puntos fijos cuando
  proceda y conserve código, datos, estilo y metadatos.
\end{enumerate}

\begin{evidencebox}{Criterio de cierre del banco}
Una actividad está terminada cuando otra persona puede repetirla, la predicción
algebraica se contrasta con los datos, el rasgo interpretado sobrevive al control
elegido y la conclusión se mantiene dentro del alcance del método. Obtener una imagen aporta
criterio de cierre.
\end{evidencebox}

\begin{warningbox}
Las actividades que requieren cálculos externos parten de datos exportados por
Toolbox Chaos y deben conservar el script o procedimiento usado. Las atribuciones deben corresponder
a la interfaz estimadores que no implementa, ni ocultar pasos manuales bajo frases
como ``el software demostró''. Cada herramienta se acredita por la operación que
realmente realizó.
\end{warningbox}

\lecturaseccion{Para seleccionar problemas adicionales: Robert L. Devaney,
\emph{A First Course in Chaotic Dynamical Systems: Theory and Experiment},
2.\textsuperscript{a} ed., capítulos 3--14; Kathleen T. Alligood, Tim D. Sauer y
James A. Yorke, \emph{Chaos: An Introduction to Dynamical Systems}, capítulos 1--13
y apéndice B; Julien C. Sprott, \emph{Elegant Chaos: Algebraically Simple Chaotic
Flows}, capítulos 1--10; Julien C. Sprott, \emph{Strange Attractors: Creating
Patterns in Chaos}, secciones 2.3--2.4, 3.1--3.5, capítulos 4--6 y apéndices E--F.}
